\documentclass{syllabus}

\begin{document}

\thispagestyle{fancy}
\date{Fall 2025}
\author{}
\title{PHY 3960: Portfolio Sequence B \\ {\large \mdseries Professional development as a physics major}}
\maketitle



\section*{General Information}

\raisebox{0.9\height}{

\begin{tabular}{lll}
\textbf{Instructor:}    & \multicolumn{2}{l}{Joseph Anderson}                             \\
\textbf{}               & \multicolumn{2}{l}{\href{mailto:janderson19@carthage.edu}{janderson19@carthage.edu}}                    \\
\textbf{}               & \multicolumn{2}{l}{DSC 082}                    \\
\textbf{}               & \multicolumn{2}{l}{(262)551-5871}                               \\ \\
\textbf{Office hours:}  & Wed, Fri               & 10:20 AM - 12:00 PM                    \\
\textbf{}               & Tuesday                & 3:30-5:00 PM (Discussion)              \\ \\
\textbf{Textbook:}      & \multicolumn{2}{l}{N/A} \\ \\
\textbf{Prerequisites:} & \multicolumn{2}{l}{PHY 2970 (Port. A)}                      
\end{tabular}
%
}

\section*{Course overview}

The second course in the Physics Portfolio Seminar sequence is focused on the personal and professional development of third year physics majors. This includes introducing students to possible career and post-baccalaureate opportunities, communicating physics ideas to nonspecialists, and encouraging them to reflect upon their own growth and accomplishments. Students in this seminar will begin in earnest the process of developing their own professional portfolio of work in physics.

\subsection*{Student learning outcomes} 
This course is structured to give students space to develop into better people primarily in the following two traits:
\begin{enumerate}[label=\textcolor{CarthageDark}{\bfseries\arabic*.}]
    \item\textbf{Magnanimity:} Aspiring to great deeds in a way that balances lofty goals and realism about personal capabilities.
    \item\textbf{Diligence:} Pursuing consistently the personal change and improvement necessary to achieve great deeds regardless of external incentives.
\end{enumerate}
Although the course will be organized so as to encourage this development, you are the only one who can determine whether to do the developing or not. In light of this, course assessment will be primarily structured around the following outcomes. 

Upon completion of this course, a student who engages thoughtfully and earnestly with the course activities will be able to:
\begin{enumerate}[label=\textcolor{CarthageDark}{\bfseries\arabic*.}]
    \item Identify societal goals worth devoting your life toward and explain how physics as a discipline can contribute meaningfully toward those goals.
    \item Explain pitfalls to professional development posed by deficiency or excess of idealistic motivation in career goals.
    \item Be aware of traditional and nontraditional physics career paths that contribute toward these goals
    \item Identify principles of seeking out knowledge and skills and how to avoid tendencies toward distractions in this process.
    \item Accurately assess the state of their knowledge and skills with respect to potential career paths and create a plan for how to develop knowledge and skill commensurate to career goals.
    \item Communicate the value of physics to a variety of audiences.
\end{enumerate}



\section*{Course components}

\subsection*{Participation} This course will involve frequent class discussions as well as student presentations; students will receive a participation score for each class session. A portion of this grade awarded based on attendance, with tardiness (arriving after 8:00) also being penalized. The remainder of the participation grade is awarded based on engagement with class activities (discussions or otherwise).

\subsection{Reflection}
This course is primarily designed as a reflective exercise for you, the student, to analyze your own motivations, goals, and immediate plan for personal development in the context of studies toward a physics degree. Reflection prompts are meant to connect class discussions to your personal experience. 

\textbf{AI usage is strictly prohibited on reflection assignments.} Reflections will be hand-written and require signing the following statement on each assignment:
\begin{quote}
    I acknowledge that the use of generative AI on this assignment represents a cowardly avoidance of engagement with my own personal experience.
\end{quote}

\subsection{Professional development plan} Over the course of the semester, there will be several assignments working toward a professional development plan. This will involve 
\begin{enumerate}
    \item Identification of possible career trajectories and your motivations for pursuing these careers.
    \item Identification of skills or experience needed to succeed in the identified career trajectories.
    \item Proposal of a timeline of concrete activities on and off campus for developing these skills or obtaining these experiences over the remainder of their undergraduate career.
    \item Development of a resume for use in applying for identified career experiences for summer 2026.
\end{enumerate}
Students are encouraged to have frequent discussions with their advisor or other mentors on this topic.



\subsection{Physics demo presentation} 
Being able to explain the value of your physics studies to a popular audience will help you to understand your motivation towards future careers. As such, in the later portion of this semester the physics department will run a booth at the Kenosha farmers market where students will bring a physics demonstration to explain both a physics principle and how this principle shows the impact of physics on society. A practice presentation will be given in class early in the semester.

\section*{Grading Scheme}
\subsection*{Component Weights}

\noindent\begin{tabular}{ll}
\textbf{Component} & \textbf{Weight}       \\ \hline
\textbf{Participation}            & 15\%  \\
\textbf{Reflection}            & 20\%       \\
\textbf{Prof. dev. plan}            &     35\%   \\
\textbf{Demo presentation}            & 30\%       \\
\end{tabular}%



\lettergrades{a) authentic engagement with class discussions and reflection assignments, and b) pattern of seeking help in developing a professional development plan from the instructor, academic advisor, or other identified mentor.}

\section*{Policies}
\subsection{General course policies}
Expect a response within one business day for emails sent to the instructor. To ensure a prompt response, format the subject line as: [PHY 3960] *insert subject here*

Students are expected to arrive and leave on time. Cf. participation section for more information. I will attempt to arrive five minutes early and be available for at least five minutes after class.

\collegepolicies[College policies]

\end{document}
